\section{构建真正衡量调查的基准}
\label{sec:design}

\textbf{RQ4} 问基准构建者能做什么。每类捷径都对应一项检查，其成本与构建基准相比很低（\cref{tab:checklist}）。下面说明其中最不显然的两项。

\begin{table}[t]
  \caption{面向基准构建者的检查项，以及每项针对的捷径类型。}
  \label{tab:checklist}
  \centering\footnotesize
  \begin{tabular}{@{}>{\raggedright\arraybackslash}p{0.75cm}>{\raggedright\arraybackslash}p{7.05cm}@{}}
    \toprule
    类型 & 检查 \\
    \midrule
    全部 & \textbf{报告不调查的下限。}在模型成绩旁边公布一条不含 LLM 的固定规则；低于它的分数不能作为调查能力的证据。 \\
    S1 & \textbf{测量目标泄漏。}按与题目的纯文本重叠对候选证据排序；如果目标排第一的频率远高于随机，就改用角色或时间来描述它，而不是用名字。 \\
    S2 & \textbf{按标签来源切分。}留出组织（或分析员），并报告仅凭来源能解释多少标签。 \\
    S3、S4 & \textbf{为每个攻击配一个良性孪生}，使用相同的工具、账户或规则，使得没有任何单条观测能区分二者。 \\
    全部 & \textbf{检查难度是否随深度变化。}按所需调查量（跳数、捷径标记）报告成功率；衡量调查的基准，在需要更多调查的地方应当成功率更低。 \\
    全部 & \textbf{公开逐题标记和日志。}标记让他人能在无捷径子集上计分；日志使 \cref{sec:scores} 这样的分析成为可能。 \\
    \bottomrule
  \end{tabular}
\end{table}

\paragraph{良性孪生消除单条观测捷径} 我们的第三个模拟环境 comp-triage 正是针对捷径 S4 构建的。它的三个场景（远程创建服务、计划任务、大量外传）各把一个攻击与三个使用相同工具的合法活动配对；以远程创建服务为例，用窃取的管理员凭据横向移动，与管理员维护、授权渗透测试和软件部署代理相互竞争。没有任何原因拥有只属于它的观测，72 个案例中有 18 个是攻击。在另外两个环境中靠"等到一条确认观测"就能解题的无 LLM 控制器，在这里只在 \compN{} 个案例中识别出 \compSingle{} 个正确原因，并把 \compSingleEsc{} 个升级。综合多条观测的控制器识别出 \compJoint{} 个，但把 \compJointMissed{} 个攻击误结为其良性孪生。配对设计让分诊重新变成必须综合证据的问题，而且综合错误的代价是漏掉攻击，这正是调查型基准应当衡量的。

\paragraph{真实日志中存在良性孪生，但公开数据集中很少} OTRF 中的伪装任务（\cref{sec:otrf}）就是天然的孪生：它的名字和目录都与 Windows 任务一致，只有创建账户和登录类型能把它区分出来。在我们审计的录制中，这是唯一的此类案例；良性事件中从来没有管理员合法地做同样的事。构建孪生需要在录制攻击的同时，用相同工具录下良性的管理操作，而我们审计的公开数据集都没有这样做。

\paragraph{抗捷径子集} 对 ExCyTIn，我们为审计过的每个题集公开逐题标记（点名、相邻、查找能否答对），以及 \lgResN{} 道抗捷径测试题的列表。在这些题上，强智能体的排名依然成立（\cref{sec:scores}），而查找按构造得零分，所以模型在这里的分数不会因点名而虚高。

\smallskip\noindent\textbf{发现 3（RQ4）。}每类捷径都有一项低成本的检查。能改变基准所衡量内容的，是报告不调查的下限、良性孪生，以及检验成功率是否随所需调查量增加而下降。
